\section{Recover BST (Two Nodes Swapped)}
\label{sec:bst-recover}

\problemstatement{The values of exactly two nodes in a BST were swapped
by mistake, breaking the BST invariant. Without changing the tree's
structure, recover the original BST by fixing the values of those two
nodes (find them and swap their values back).}

\begin{patternbox}
\emph{``Exactly two nodes are out of place''} in a problem framed around
\textbf{sortedness} is the strongest possible signal to invoke
Section~\ref{sec:bst-intro}'s headline fact directly: a correct BST's
inorder traversal is strictly increasing, so the swapped pair must show
up as exactly one or two places where the inorder sequence decreases
instead of increases -- find those violation points during a single
inorder traversal, and the swapped nodes are sitting right there.
\end{patternbox}

\subsection*{Understanding the problem carefully}

In a correct BST, scanning the inorder sequence left to right, every
value is greater than the one before it. With two values swapped, this
property breaks at either one place (if the two swapped nodes happen to
be \emph{adjacent} in inorder order) or two places (if they are not
adjacent) -- and at each violation, the \emph{earlier} of the two
compared values is the one that is too large for its position, or the
\emph{later} one is too small.

\begin{ideabox}[Track the First and Second Violation During One Inorder Pass]
Maintain $\textit{prev}$ (the previously visited node) and two
not-yet-found pointers, $\textit{first}$ and $\textit{second}$. During
inorder traversal, at each node: if $\textit{prev}$ exists and
$\textit{prev}.\textit{val} > \textit{node}.\textit{val}$ (a violation):
if $\textit{first}$ has not yet been set, set $\textit{first} \gets
\textit{prev}$; in either case, set $\textit{second} \gets
\textit{node}$ (always update to the most recent violation's second
node). Update $\textit{prev} \gets \textit{node}$ and continue. After
the traversal, swap $\textit{first}.\textit{val}$ and
$\textit{second}.\textit{val}$.
\end{ideabox}

\subsection*{Why exactly these two recorded nodes are the swapped pair}

If the swap is between adjacent inorder values, exactly one violation
occurs, and $\textit{first}$/$\textit{second}$ from that single violation
are exactly the two swapped nodes. If the swap is between non-adjacent
values, exactly \emph{two} violations occur: the first violation's
\textit{prev} is the node that ended up holding too large a value (it
should be smaller, since the swap displaced a large value into an
earlier position), and the second violation's current node is the one
holding too small a value (displaced later than it should be). Always
keeping $\textit{first}$ fixed at the \emph{first} violation encountered,
while letting $\textit{second}$ update to track the \emph{most recent}
violation's later node, correctly identifies both endpoints of the
original swap regardless of which of these two cases occurred --
swapping their values back is enough to restore strict monotonicity,
since no third node is ever out of place.

\subsection*{Algorithm}

\begin{enumerate}
  \item Initialize $\textit{prev}, \textit{first}, \textit{second}
        \gets$ null.
  \item Perform an inorder traversal. At each visited node:
  \begin{enumerate}
    \item If $\textit{prev} \ne$ null and $\textit{prev}.\textit{val} >
          \textit{node}.\textit{val}$:
    \begin{itemize}
      \item If $\textit{first} = $ null: $\textit{first} \gets
            \textit{prev}$.
      \item $\textit{second} \gets \textit{node}$.
    \end{itemize}
    \item $\textit{prev} \gets \textit{node}$.
  \end{enumerate}
  \item Swap $\textit{first}.\textit{val}$ and
        $\textit{second}.\textit{val}$.
\end{enumerate}

\subsection*{Dry run}

\dryrun{A BST whose intended inorder sequence is $1,2,3,4,5,6,7$ has
nodes holding values $3$ and $6$ swapped, giving the actual inorder
sequence $1,2,6,4,5,3,7$.}

\begin{itemize}
  \item Visit $1$: $\textit{prev}=$ null, no check. $\textit{prev}=1$.
  \item Visit $2$: $1>2$? No. $\textit{prev}=2$.
  \item Visit $6$: $2>6$? No. $\textit{prev}=6$.
  \item Visit $4$: $6>4$? Yes -- violation. $\textit{first}=$ null, so
        $\textit{first} \gets$ node holding $6$. $\textit{second} \gets$
        node holding $4$. $\textit{prev}=4$.
  \item Visit $5$: $4>5$? No. $\textit{prev}=5$.
  \item Visit $3$: $5>3$? Yes -- violation. $\textit{first}$ already
        set, unchanged. $\textit{second} \gets$ node holding $3$.
        $\textit{prev}=3$.
  \item Visit $7$: $3>7$? No. $\textit{prev}=7$.
\end{itemize}

$\textit{first}$ is the node currently holding $6$; $\textit{second}$ is
the node currently holding $3$. Swap their values: the node that held $6$
now holds $3$, and the node that held $3$ now holds $6$. The inorder
sequence becomes $1,2,3,4,5,6,7$ -- correctly recovered.

\subsection*{C++ implementation}

\begin{lstlisting}
#include <bits/stdc++.h>
using namespace std;

struct TreeNode {
    int val;
    TreeNode* left;
    TreeNode* right;
    TreeNode(int v) : val(v), left(nullptr), right(nullptr) {}
};

void inorder(TreeNode* node, TreeNode*& prev, TreeNode*& first, TreeNode*& second) {
    if (!node) return;

    inorder(node->left, prev, first, second);

    if (prev && prev->val > node->val) {
        if (!first) first = prev;
        second = node;
    }
    prev = node;

    inorder(node->right, prev, first, second);
}

void recoverTree(TreeNode* root) {
    TreeNode *prev = nullptr, *first = nullptr, *second = nullptr;
    inorder(root, prev, first, second);
    swap(first->val, second->val);
}

int main() {
    // Built so that an inorder traversal visits 1, 2, 6, 4, 5, 3, 7 --
    // i.e. the values 3 and 6 have been swapped relative to a correct BST.
    TreeNode* root = new TreeNode(4);
    root->left = new TreeNode(2);
    root->left->left = new TreeNode(1);
    root->left->right = new TreeNode(6); // swapped value
    root->right = new TreeNode(3);       // swapped value
    root->right->left = new TreeNode(5);
    root->right->right = new TreeNode(7);

    recoverTree(root);

    TreeNode *p = nullptr, *f = nullptr, *s = nullptr;
    inorder(root, p, f, s); // verify: no violations should remain
    cout << (f == nullptr ? "Recovered" : "Still broken") << endl;
    return 0;
}
\end{lstlisting}

\begin{complexitybox}
\textbf{Time Complexity.} $O(n)$ -- a single inorder traversal visiting
every node once.
\textbf{Space Complexity.} $O(h)$ for the recursion stack. (An advanced
variant using Morris traversal achieves $O(1)$ space, at the cost of
temporarily modifying the tree's pointers during the scan -- worth
knowing exists, though the recursive version here is the one to derive
first.)
\end{complexitybox}

\begin{takeawaybox}
``Exactly $k$ elements are out of place in an otherwise sorted sequence''
problems are solved by scanning once and recording violation points,
rather than sorting from scratch. Here, $k=2$ gives exactly the
\textit{first}/\textit{second} bookkeeping shown; the same idea
generalizes (with more bookkeeping) to other small, fixed values of $k$.
\end{takeawaybox}
